\documentclass[11pt]{article}
\usepackage[margin=1in]{geometry}
\usepackage{amsmath,amssymb,amsthm}
\usepackage{graphicx}
\usepackage[hidelinks]{hyperref}

\newtheorem{theorem}{Theorem}
\newtheorem{corollary}{Corollary}
\newtheorem{remark}{Remark}

\newcommand{\K}{\mathcal K}
\newcommand{\B}{\mathcal B}
\newcommand{\U}{\mathcal U}
\newcommand{\wh}{\widehat}
\newcommand{\eps}{\varepsilon}

\title{The Exact Noncompactness Norm of the Calkin Representation}
\author{FA Banach run \texttt{fa\_banach\_001}, agent \texttt{agent\_lane\_05}}
\date{2026-06-16}

\begin{document}
\maketitle

\section*{Status}

This packet records a candidate partial result, likely valid.  It sharpens the
previous lane-5 packet on Remark 1 after Theorem 2.1 of Boedihardjo--Johnson,
arXiv:1301.0054.  The source question asks whether their Calkin
representation theorem remains true without assuming that \(X\) has the
bounded compact approximation property (BCAP).

\section*{Source Context}

Boedihardjo and Johnson define, for each infinite-dimensional Banach space
\(Y\), the ultrapower space
\[
  \widehat Y=\{(y_\alpha^*)_{\alpha,\U}\in (Y^*)^\U:
  w^*\hbox{-}\lim_{\alpha,\U}y_\alpha^*=0\},
\]
where \(\alpha\) ranges over finite-dimensional subspaces of a fixed
infinite-dimensional Banach space \(X\).  If \(T\in B(X,Y)\), they define
\(\widehat T\in B(\widehat Y,\widehat X)\) by
\[
  (y_\alpha^*)_{\alpha,\U}\mapsto (T^*y_\alpha^*)_{\alpha,\U}.
\]
They observe that \(\widehat K=0\) for compact \(K\).  Their Theorem 2.1
states that, if \(X\) has the \(\lambda\)-BCAP, then
\[
  f:B(X)/K(X)\to B(\widehat X),\qquad \dot T\mapsto\widehat T,
\]
is a norm-one \((\lambda+1)\)-embedding.

\begin{center}
\includegraphics[width=0.98\linewidth]{figures/theorem_context_crop.png}
\end{center}

Immediately after the proof, the authors ask whether Theorem 2.1 remains true
without the BCAP hypothesis.

\begin{center}
\includegraphics[width=0.98\linewidth]{figures/open_problem_crop.png}
\end{center}

\section*{Result}

For \(T\in B(X)\) and a finite-dimensional subspace \(E\subset X\), write
\(Q_E:X\to X/E\) for the quotient map and set
\[
  \kappa_X(T):=\inf_{E\subset X,\ \dim E<\infty}\|Q_E T\|.
\]
This is the limiting Kolmogorov noncompactness quantity of the set \(T(B_X)\).

\begin{theorem}
Let \(X\) be an infinite-dimensional Banach space and use the
Boedihardjo--Johnson ultrafilter on the directed set of finite-dimensional
subspaces of \(X\).  Then, for every \(T\in B(X)\),
\[
  \|\widehat T\|=\kappa_X(T)
  =\inf_{E\subset X,\ \dim E<\infty}\|Q_E T\|.
\]
\end{theorem}

\section*{Proof Intuition}

The lower bound is the construction from the previous packet: for each
finite-dimensional \(E\), take a norm-one functional in \(E^\perp\) on which
\(T^*\) almost attains \(\|T^*|_{E^\perp}\|=\|Q_E T\|\).  These functionals are
weak-star-null along the tail ultrafilter, so they give a test vector in
\(\widehat X\).

The new point is the reverse inequality.  Any vector of \(\widehat X\) is
represented by a bounded weak-star-null family \((x_\alpha^*)\).  Fix a
finite-dimensional \(F\subset X\).  Since weak-star convergence on \(F\) is
norm convergence in \(F^*\), the restrictions \(x_\alpha^*|_F\) are small on
an ultrafilter-large set.  Equivalently, \(x_\alpha^*\) is close to
\(F^\perp\).  Replacing \(x_\alpha^*\) by a nearby element of \(F^\perp\)
forces \(T^*x_\alpha^*\) to be bounded by \(\|T^*|_{F^\perp}\|=\|Q_FT\|\).
Taking the infimum over \(F\) gives the upper bound.

\begin{proof}
Let \(\Lambda_0\) be the directed set of finite-dimensional subspaces of \(X\),
ordered by inclusion, and let \(\U\) be an ultrafilter containing every tail
\[
  \{\alpha\in\Lambda_0:\alpha\supset \alpha_0\}.
\]
For \(E\in\Lambda_0\), duality gives
\[
  \|Q_E T\|=\|(Q_E T)^*\|=\|T^*|_{E^\perp}\|.
\]
The function \(E\mapsto\|Q_E T\|\) is decreasing on \(\Lambda_0\), so its
\(\U\)-limit is its infimum, namely \(\kappa_X(T)\).

We first prove \(\|\widehat T\|\ge\kappa_X(T)\).  Fix \(\eps>0\).  For every
\(E\in\Lambda_0\), choose \(x_E^*\in E^\perp\) with \(\|x_E^*\|=1\) and
\[
  \|T^*x_E^*\|\ge \|T^*|_{E^\perp}\|-\eps
  =\|Q_E T\|-\eps.
\]
The family \((x_E^*)_{E,\U}\) belongs to \(\widehat X\), because for each
\(x\in X\), the tail \(E\supset \operatorname{span}\{x\}\) lies in \(\U\), and
on that tail \(x_E^*(x)=0\).  Therefore
\[
  \|\widehat T\|
  \ge \lim_{E,\U}\|T^*x_E^*\|
  \ge \lim_{E,\U}\|Q_E T\|-\eps
  =\kappa_X(T)-\eps.
\]
Letting \(\eps\downarrow0\) gives the lower bound.

For the upper bound, let
\(\xi=(x_E^*)_{E,\U}\in\widehat X\) with \(\|\xi\|\le1\).  Fix a
finite-dimensional subspace \(F\subset X\) and \(\eps>0\).  Since
\(w^*\hbox{-}\lim_{E,\U}x_E^*=0\), and since \(F\) is finite dimensional,
\[
  \lim_{E,\U}\|x_E^*|_F\|_{F^*}=0.
\]
For \(E\) in an ultrafilter-large set, choose \(z_E^*\in F^\perp\) with
\[
  \|x_E^*-z_E^*\|\le \|x_E^*|_F\|+\eps.
\]
This is possible because the quotient \(X^*/F^\perp\) is isometric to
\(F^*\) by restriction.  Hence, on an ultrafilter-large set,
\(\|x_E^*-z_E^*\|\le 2\eps\).  It follows that
\[
  \|T^*x_E^*\|
  \le \|T^*z_E^*\|+2\eps\|T\|
  \le \|T^*|_{F^\perp}\|\|z_E^*\|+2\eps\|T\|.
\]
Also \(\|z_E^*\|\le \|x_E^*\|+2\eps\).  Taking \(\U\)-limits and using
\(\lim_{E,\U}\|x_E^*\|=\|\xi\|\le1\), we obtain
\[
  \|\widehat T\xi\|
  \le \|T^*|_{F^\perp}\|(1+2\eps)+2\eps\|T\|
  =\|Q_FT\|(1+2\eps)+2\eps\|T\|.
\]
Letting \(\eps\downarrow0\), and then taking the infimum over
finite-dimensional \(F\subset X\), gives
\[
  \|\widehat T\xi\|\le \kappa_X(T).
\]
Since this holds for every \(\xi\in\widehat X\) with \(\|\xi\|\le1\), we have
\(\|\widehat T\|\le\kappa_X(T)\).  This completes the proof.
\end{proof}

\begin{corollary}
The Boedihardjo--Johnson representation is injective modulo compact operators
for every infinite-dimensional Banach space \(X\).
\end{corollary}

\begin{proof}
The quantity \(\kappa_X(T)\) is zero if and only if \(T\) is compact.  Indeed,
if \(\kappa_X(T)=0\), then for every \(\eps>0\), the set \(T(B_X)\) lies within
\(\eps\) of a finite-dimensional subspace of \(X\).  Since \(T(B_X)\) is
bounded, this makes \(T(B_X)\) totally bounded.  Conversely compact sets are
approximated by finite-dimensional spans of finite nets.  The theorem now
gives \(\widehat T=0\) if and only if \(T\) is compact.
\end{proof}

\section*{Reduction of the Remaining Open Problem}

The theorem identifies exactly what the ultrapower representation measures.
Thus Theorem 2.1 without BCAP is equivalent to asking whether, for the relevant
Banach space \(X\), there is a constant \(C_X<\infty\) such that
\[
  \|T\|_e\le C_X\,\kappa_X(T)
  \qquad (T\in B(X)).
\]
The BCAP proof in Boedihardjo--Johnson proves precisely this estimate with
\(C_X=\lambda+1\) when \(X\) has the \(\lambda\)-BCAP.  Without BCAP, the
representation remains faithful, but the unresolved issue is whether the
essential norm and the Kolmogorov noncompactness norm are uniformly comparable
on \(B(X)/K(X)\).

Remark 2 is similarly reframed: for \(X=L_p(0,1)\), \(p\ne2\), the
representation is isometric exactly when
\[
  \|T\|_e=\kappa_X(T)
  \qquad (T\in B(X)).
\]

\section*{Novelty and Literature Check}

The run indexes were searched for arXiv:1301.0054, Boedihardjo--Johnson,
Calkin, BCAP, essential norm, and measure/Kolmogorov noncompactness phrases.
The only direct run hit was the prior lane-5 injectivity packet.  Bounded web
searches for exact phrases around Theorem 2.1, the BCAP-free Calkin
representation, and essential norm versus measure of noncompactness did not
surface a later paper resolving Remark 1.  Novelty confidence is moderate: the
norm formula is natural once one looks at the quotient norms, but it is not
recorded in the run memory and was not visible in the bounded search.

\section*{Limitations}

This packet does not settle whether the estimate
\(\|T\|_e\le C_X\kappa_X(T)\) holds for every Banach space, nor does it produce
a space and operators for which the ratio \(\|T\|_e/\kappa_X(T)\) is
unbounded.  It should therefore be treated as a sharp reduction and a genuine
strengthening of the injectivity result, not as a full solution of Remark 1.

\section*{Human Review Recommendation}

Reviewers should check the upper-bound argument in the theorem, especially the
use of finite-dimensionality of \(F\) to turn weak-star convergence into norm
convergence of restrictions to \(F\), and the isometric identification
\(X^*/F^\perp\simeq F^*\).

\begin{thebibliography}{9}

\bibitem{BJ}
M. T. Boedihardjo and W. B. Johnson,
\emph{On Mean Ergodic Convergence in the Calkin Algebras},
arXiv:1301.0054.

\end{thebibliography}

\end{document}
